\documentclass[]{article}

\usepackage{amsmath}
\usepackage{multirow}
\usepackage{natbib}
\usepackage{graphicx} 


\begin{document}

The transport of passive particles in chaotic fluid flows is a fundamental problem in non-equilibrium statistical mechanics, with applications ranging from pollutant dispersal in the atmosphere and oceans to mixing processes in planetary interiors. In such systems, the intricate, deterministic motion of the fluid often drives particle trajectories that deviate strongly from classical Brownian motion. A common manifestation of this is anomalous superdiffusion, where the mean squared displacement grows faster than linearly with time, $\langle \Delta r^2(t) \rangle \propto t^\alpha$, with an anomalous exponent $\alpha > 1$. This behavior signals the presence of long-range correlations in the particle's velocity, posing a significant challenge to theoretical modeling: can the complex dynamics generated by a deterministic system be effectively captured by a simpler, canonical stochastic process?

The Lévy walk model has emerged as a leading theoretical framework for describing superdiffusion in physical systems. It conceptualizes a particle performing a sequence of ballistic flights, where the duration of each flight is drawn from a heavy-tailed, power-law distribution. This construction naturally gives rise to superdiffusive scaling while respecting a finite particle velocity, making it a physically plausible model for transport phenomena. The apparent success of Lévy walks in reproducing the macroscopic scaling observed in various chaotic systems raises a critical question: does this kinematic similarity imply a deeper statistical equivalence? It remains an open problem whether the full statistical structure of transport in a deterministic, chaotic Hamiltonian flow can be faithfully mapped onto a purely stochastic process like a Lévy walk.

To address this question, we conduct a direct and systematic comparison between tracer transport in a chaotic point-vortex system and a benchmark Lévy walk process. The two-dimensional point-vortex flow serves as an ideal model system; it is a conservative Hamiltonian flow that represents a foundational model for ideal 2D turbulence, and its degree of chaoticity can be precisely controlled by varying the number of vortices, $N$. We numerically simulate the trajectories of passive tracers advected by these flows for systems with increasing vortex density and compare their statistical properties against a comprehensive dataset of Lévy walk trajectories. Our analysis employs a suite of diagnostics, including the scaling of the mean squared displacement, velocity autocorrelation functions, residence time distributions in low-velocity regions, and the evolution of displacement probability distributions.

Our investigation first confirms that as the vortex density increases, the point-vortex system kinematically mimics a Lévy process. The transport transitions from near-normal diffusion to strong superdiffusion, accompanied by the emergence of power-law residence times and heavy-tailed displacement profiles, which are hallmark signatures of Lévy-like dynamics. However, the central contribution of this work is to move beyond these kinematic measures and probe the ergodicity of the underlying statistical ensembles. We uncover a fundamental divergence between the two systems. As the superdiffusive behavior in the vortex system strengthens with increasing $N$, the system becomes progressively more ergodic. This trend is in stark opposition to the behavior of Lévy walks, which become increasingly non-ergodic as their superdiffusive character intensifies. This finding reveals that despite reproducing key macroscopic signatures of a Lévy process, the deterministic and Hamiltonian constraints of the vortex flow impose a distinct statistical structure, precluding a complete equivalence with its stochastic counterpart and highlighting that kinematic similarity alone is an insufficient criterion for model validation.

\end{document}
                